% ============================================================
% 1.X  Project Approach
% ============================================================
\subsection{Project Approach}
\label{sec:intro-approach}

In the following section I discuss the various phases of the project and the approach taken at each phase.
The project followed an iterative research-driven methodology, where each phase informed and refined the next.
Any technologies, tools, or design decisions adopted during a particular phase are stated and justified.

\subsubsection{Planning Phase -- Literature Review and Feasibility}

The first phase of the project focused on understanding the problem domain and assessing whether structural homology transfer was a viable signal for PPI prediction.
I began by conducting a literature review of existing computational PPI prediction methods, including co-evolution-based approaches, machine-learning classifiers, and physics-based docking, to understand their respective strengths and limitations.
This review established that while structural similarity had been used extensively for function annotation, its systematic application to interaction prediction at proteome scale remained underexplored.

I then studied the key tools and datasets that would underpin the project: the AlphaFold Protein Structure Database~\cite{alphafold2}, the TED dataset and its domain pair list~\cite{ted_dataset}, and the Merizo-search pipeline comprising Merizo~\cite{merizo} for domain segmentation and Foldclass~\cite{merizo_search} for structural search.
Understanding the internal architecture of Merizo-search --- particularly how Foldclass encodes domains as fixed-length embeddings and retrieves candidates via cosine similarity followed by TM-align~\cite{tmalign} re-ranking --- was essential before any implementation work could begin.

During this phase I also identified the Zhang et~al.\ benchmark~\cite{zhang2025ppi} as the most appropriate evaluation framework.
Its use of weighted precision--recall to account for the extreme class imbalance in proteome-wide screening (approximately 1:1{,}000 signal-to-noise ratio) made it particularly well suited to assessing the practical utility of the approach, rather than reporting artificially inflated metrics on a balanced test set.

\subsubsection{Design Phase -- Pipeline Architecture and Database Construction}

With the feasibility established, the second phase focused on designing the end-to-end pipeline architecture.
The central design challenge was transforming the TED domain pair list --- which catalogues over 200~million domain entries with intra-protein domain--domain contacts --- into a Foldclass-compatible indexed database that could be searched efficiently.

I designed the pipeline around a five-stage workflow: (1)~query protein segmentation into domains using Merizo, (2)~lookup of known domain--domain contacts from the TED pair list to identify interface domains, (3)~extraction of the interface domain (Domain~$B$) as the search query, (4)~Foldclass embedding search over the indexed database to retrieve structurally similar candidates, and (5)~TM-align re-ranking and scoring of the top candidates to produce a ranked list of predicted interaction partners.

To support targeted biological queries, I designed a SQLite metadata index storing per-domain information including taxonomy~ID, CATH fold classification, AlphaFold model confidence (pLDDT), and domain density.
This schema was designed to support query-time filtering, so that a researcher could restrict a search to, for example, all human kinase domains without loading the entire 200-million-domain database into memory.
A filtered embedding iterator was designed alongside the metadata index, loading only the relevant subset of precomputed embeddings during search to reduce both runtime and memory consumption.

\subsubsection{Execution Phase -- Implementation}

The implementation followed an incremental approach.
I began with the database construction scripts, converting the raw TED pair list into the binary embedding format and SQLite index required by Foldclass.
Once the database was in place, I implemented the core search pipeline, initially testing it on small subsets of the data to verify correctness before scaling to the full database.

The filtered embedding iterator was implemented next, as it was required for practical operation on the full-scale database.
Without filtering, each search would load all 200~million domain embeddings into memory, which was infeasible on available hardware.
The iterator reads the SQLite index first, determines which embedding files contain domains matching the query filters, and loads only those files.

The end-to-end workflow was wrapped in a shell script (\texttt{search-potential-interactions.sh}) that orchestrates the full pipeline from an input protein structure to a ranked list of candidate interaction partners.
Version control via Git was used throughout to track changes, and regular meetings with my supervisor ensured the implementation remained aligned with the project goals.

\subsubsection{Validation Phase -- Benchmarking and Analysis}

The final phase focused on systematic evaluation against the Zhang et~al.\ benchmark~\cite{zhang2025ppi}.
I implemented a benchmarking script (\texttt{benchmark\_ppi.py}) that runs the pipeline over the 3{,}000 positive and 30{,}000 negative pairs in the benchmark, computes TM-scores for each pair, and generates weighted precision--recall curves.

A sweep of TM-score thresholds from 0.3 to 0.9 was conducted to characterise the trade-off between precision and recall.
At each threshold, I computed area under the weighted precision--recall curve (AUCPR) and precision at recall thresholds of 0.1, 0.2, and 0.5.
Beyond aggregate metrics, I performed a qualitative failure analysis, manually inspecting cases where the method produced false positives or false negatives to identify the principal failure modes --- including paralogs with divergent interaction partners, convergent fold evolution, and interactions mediated by intrinsically disordered regions that lack stable tertiary structure.

The filtering infrastructure was also validated by comparing runtime and peak memory usage with and without SQLite-based filtering, to quantify the practical gains of the metadata index and filtered iterator.


% ============================================================
% 1.X  Report Structure
% ============================================================
\subsection{Report Structure}
\label{sec:intro-report-structure}

This report documents the design, implementation, and evaluation of a proteome-scale PPI prediction pipeline based on domain structural homology transfer.
Below is an outline of the structure this report will follow:

\begin{itemize}
    \item \textbf{Chapter~\ref{ch:background}: Background and Related Work} ---
    A review of the biological context of protein--protein interactions, existing computational PPI prediction methods, and the enabling technologies underpinning this work: AlphaFold2, the TED dataset, CATH domain classification, and the Merizo-search pipeline.
    The chapter concludes with a discussion of structural homology transfer and the evaluation framework adopted.

    \item \textbf{Chapter~\ref{ch:design}: System Design and Implementation} ---
    A detailed description of the pipeline architecture, including database construction from the TED pair list, the search pipeline design, the SQLite metadata index schema, the filtered embedding iterator, and the end-to-end workflow.
    Key design decisions and the alternatives considered are discussed.

    \item \textbf{Chapter~\ref{ch:evaluation}: Evaluation Methodology} ---
    A description of the Zhang et~al.\ benchmark, the weighted precision--recall framework, the metrics used, and the TM-score threshold sweep design.

    \item \textbf{Chapter~\ref{ch:results}: Results and Analysis} ---
    Presentation and analysis of the benchmark results, including threshold sweep curves, case studies of successful and failed predictions, a critical failure analysis, and validation of the filtering infrastructure.

    \item \textbf{Chapter~\ref{ch:conclusions}: Conclusions} ---
    A summary of the project's contributions, a critical evaluation of the approach, discussion of limitations and future work, and reflections on the project.
\end{itemize}
